\subsection{The non-self-dual quadratic family}

For an integer $m\geq0$, let $S(m)$ be the commutative based ring with basis
$\{1,g,x,y\}$, involution $g^*=g$, $x^*=y$, and multiplication
\[
 g^2=1,\qquad gx=y,\qquad gy=x,\qquad
 x^2=y^2=g+m(x+y),\qquad xy=1+m(x+y).
\]

\begin{theorem}\label{thm:nonselfdual-quadratic}
For every $m\geq2$, the fusion ring $S(m)$ admits no fusion categorification
over $\mathbb C$.
\end{theorem}

\begin{proof}
Put $d=m+\sqrt{m^2+1}$. Thus $d^2=1+2md$, and the Frobenius--Perron
dimensions are $(1,1,d,d)$. By Schopieray's Galois-conjugacy theorem for
two-dimension fusion rings \cite[Lemma~3.1]{Schopieray}, any fusion
categorification would have a pseudo-unitary Galois conjugate with the same
fusion rules. It is consequently sufficient to exclude pseudo-unitary
categorifications.

Suppose one exists, and use its canonical spherical structure. Its dimension is
$D=2+2d^2=4+4md$. The four characters of $S(m)$, displayed by columns, are
\[
 \begin{pmatrix}
 1&1&1&1\\
 1&1&-1&-1\\
 d&-d^{-1}&i&-i\\
 d&-d^{-1}&-i&i
 \end{pmatrix}.
\]
Their formal codegrees are $D$, $2+2d^{-2}$, $4$, $4$, respectively.
The formal-codegree description of induction therefore gives
\[
 I(\mathbf1)\cong\mathbf1\oplus A\oplus B\oplus C,
 \qquad
 \operatorname{FPdim}(A)=d^2=1+2md,
 \quad
 \operatorname{FPdim}(B)=\operatorname{FPdim}(C)=1+md,
\]
where the four central simples are pairwise nonisomorphic and occur with
multiplicity one. Indeed, every irreducible representation of this commutative
rank-four ring is one-dimensional. Furthermore,
\[
 F I(\mathbf1)\cong
 \bigoplus_{T\in\operatorname{Irr}\mathcal C}T\otimes T^*
 \cong 4\mathbf1\oplus2mX\oplus2mY,
\]
where $X^*\cong Y$ correspond to $x,y$. By adjunction each of $F(A)$,
$F(B)$, $F(C)$ contains $\mathbf1$ once, and the displayed decomposition
shows that none contains $g$. Since $d$ is irrational, comparison of dimensions
now forces
\[
 [F(A):X]+[F(A):Y]=2m,\qquad
 [F(B):X]+[F(B):Y]=[F(C):X]+[F(C):Y]=m.
\]
These statements determine the sum of the two low-part dimensions without
requiring the individual restriction multiplicities:
\begin{align*}
 L&=\operatorname{FPdim}(I_{\mathrm{low}}(X))+
    \operatorname{FPdim}(I_{\mathrm{low}}(Y))\\
  &=2m(1+2md)+2m(1+md)=4m+6m^2d.
\end{align*}
Both $X$ and $Y$ are non-self-dual. Applying
\eqref{eq:nonselfdual-third} to each of them and adding gives $3L\leq2dD$.
But
\begin{align*}
 3L-2dD
 &=12m+18m^2d-2d(4+4md)\\
 &=12m+18m^2d-8d-8m(1+2md)\\
 &=4m+(2m^2-8)d>0
\end{align*}
for $m\geq2$, a contradiction.
\end{proof}

